\section{Related Work}
\label{sec:related}

\subsection{Dark-Vessel Detection in SAR}

SAR vessel detection combines small-object localization with a sensing domain
whose image formation differs substantially from natural imagery. Speckle,
incidence-angle effects, sea state, sidelobes, and bright coastal
infrastructure can all alter the appearance of a vessel-sized response. The
xView3-SAR benchmark couples Sentinel-1 imagery with vessel and fixed-object
annotations and evaluates geographically matched detections rather than box
overlap~\cite{paolo2022xview3}. Its dark-vessel objective also distinguishes
the operational question from generic ship detection: the relevant final
slice consists of human-verified vessels lacking a matched AIS report. We use
the benchmark's point geometry directly and reserve the human-verified scenes
for a single final evaluation.

Many SAR ship detectors are built around bounding boxes. Boxes are useful when
extent is reliably annotated, but deriving them from point labels would add a
size model that is unnecessary for the xView3 distance metric. We instead use
a CenterNet-style representation in which objects are points on a heatmap and
are decoded by local peak extraction~\cite{zhou2019objects}. TRANSAR provides
recent SAR precedent for heatmap-based object localization and SAR-aware
representation learning~\cite{almalioglu2025transar}. Our detector is not
presented as a methodological contribution; it is held fixed to expose the
effect of encoder initialization.

\subsection{Generic and Remote-Sensing Pretraining}

ImageNet pretraining is a common baseline for visual transfer
learning~\cite{deng2009imagenet}. In the ViT track, the generic control is a
ViT-B/16 checkpoint trained with supervised augmentation and regularization
(AugReg)~\cite{dosovitskiy2021image,steiner2021trainvit}. In the CNN track, the
generic control is ConvNeXt-V2-Base initialized by fully convolutional masked
autoencoding followed by supervised ImageNet-1K fine-tuning
~\cite{woo2023convnextv2}. The checkpoints therefore match in final source
dataset and classification supervision, but not in full training history.
They are treated as matched \emph{roles}, not as a controlled comparison of
pretraining algorithms.

Remote-sensing pretraining may narrow the gap between source and target
statistics. SatDINO applies DINO self-distillation~\cite{caron2021emerging} to
fMoW-RGB optical overhead imagery~\cite{christie2018fmow}, yielding the optical
ViT initialization used here~\cite{straka2025satdino}. SARMAE applies masked
autoencoding~\cite{he2022masked} to a million-scale SAR corpus and supplies the
SAR ViT initialization~\cite{liu2026sarmae}. Because these released ViT
checkpoints differ in objective as well as sensing domain, their contrast is a
best-available-model comparison rather than a pure modality intervention.

The CNN track provides a tighter source-domain comparison. Its optical and SAR
arms use ConvNeXt-V2-Base checkpoints trained for the same BigEarthNet v2.0
land-cover task, using Sentinel-2 and Sentinel-1 inputs, respectively
~\cite{sumbul2019bigearthnet,clasen2025reben}. Architecture, dataset family,
and supervised objective are shared while sensing modality changes. We
therefore give the within-CNN Sentinel-1/Sentinel-2 contrast the strongest
domain-specific interpretation and use the ViT contrast as corroborating or
contrasting evidence with its method caveat visible.

\subsection{Architecture and Label Efficiency}

ViT-B/16~\cite{dosovitskiy2021image} and ConvNeXt-V2-Base
~\cite{woo2023convnextv2} are similarly sized modern encoders, but their
tokenization, spatial hierarchy, and inductive biases differ. Published
pretrained checkpoints often entangle those architectural differences with
source data and objectives. Our two-track design avoids interpreting a single
ViT--CNN score difference as a domain effect. Domain claims are made within an
architecture; cross-track comparisons pair random with random, ImageNet with
ImageNet, optical with optical, and SAR with SAR to ask whether the pattern is
architecture-general.

Label-efficiency studies further require more than reporting performance at a
single fraction. Non-nested samples can confound fraction with scene choice,
while chip-level splits can leak adjacent content across train and evaluation
sets. We use one frozen, nested sequence of scene sets and report the full
four-point curve for each arm. Because the study uses a single fixed seed, the
curves are point estimates; we do not draw uncertainty bands, estimate seed
variance, or claim statistical significance. This constraint favors careful
effect reporting---absolute F1, improvement over the within-track random
floor, and the SAR-minus-optical contrast at each registered fraction---over
post-hoc significance language.
